%% kapa: exam=1 round=35 year=2024 period=1 form=A subject=civil_law subject_ko=민법 seq=3 number=3
%% source: Q-Net 감정평가사 기출문제 (cst00309) artlSeq=5213229 file=raw/1008/5213229/1교시 문제지 원본.pdf
%% conversion: kordoc 4.24.51 -> tools/gen1.py (auto); review=proofread-sonnet; flags=none
\begin{question}{3}
의사무능력자 甲은 乙은행으로부터 5천만 원을 차용하는 대출거래약정을 체결하면서 그 담보로 자신의 X부동산에 근저당권을 설정하고 乙명의로 그 설정등기를 마쳐주었다. 이에 관한 설명으로 옳은 것을 모두 고른 것은? (다툼이 있으면 판례에 따름)
\begin{framedtext}
ㄱ. 甲과 乙이 체결한 대출거래약정 및 근저당권설정계약은 무효이다. \par
ㄴ. 甲은 그 선의·악의를 묻지 않고 乙에 대하여 현존이익을 반환할 책임이 있다. \par
ㄷ. 만약 甲이 乙로부터 대출받은 금원을 곧바로 丙에게 다시 대여하였다면, 乙은 甲에게 丙에 대한 부당이득반환채권의 양도를 구할 수 있다. \par
\end{framedtext}
\begin{choices}
  \item ㄱ
  \item ㄴ
  \item ㄷ
  \item ㄱ, ㄴ
  \item ㄱ, ㄴ, ㄷ
\end{choices}
\end{question}
\answer{5}
